% Components expose stable environments. Themes supply box styles and can
% replace these font hooks before the environments are first used.
\NewDocumentCommand{\FayorgTheoremHeadFont}{}{\bfseries}
\NewDocumentCommand{\FayorgTheoremBodyFont}{}{\itshape}
\NewDocumentCommand{\FayorgDefinitionHeadFont}{}{\bfseries}
\NewDocumentCommand{\FayorgDefinitionBodyFont}{}{\normalfont}
\NewDocumentCommand{\FayorgExampleHeadFont}{}{\bfseries}
\NewDocumentCommand{\FayorgExampleBodyFont}{}{\normalfont}
\NewDocumentCommand{\FayorgRemarkHeadFont}{}{\itshape}
\NewDocumentCommand{\FayorgRemarkBodyFont}{}{\normalfont}
\NewDocumentCommand{\FayorgExerciseHeadFont}{}{\bfseries}
\NewDocumentCommand{\FayorgExerciseBodyFont}{}{\normalfont}
\NewDocumentCommand{\FayorgExplanationHeadFont}{}{\bfseries}
\NewDocumentCommand{\FayorgExplanationBodyFont}{}{\normalfont}
\NewDocumentCommand{\FayorgAnswerHeadFont}{}{\bfseries}
\NewDocumentCommand{\FayorgAnswerBodyFont}{}{\normalfont}
\NewDocumentCommand{\FayorgConjectureHeadFont}{}{\bfseries}
\NewDocumentCommand{\FayorgConjectureBodyFont}{}{\normalfont}

\newtheoremstyle{fayorg-theorem}
  {0pt}{0pt}{\FayorgTheoremBodyFont}{}
  {\FayorgTheoremHeadFont}{.}{.5em}{}
\newtheoremstyle{fayorg-definition}
  {0pt}{0pt}{\FayorgDefinitionBodyFont}{}
  {\FayorgDefinitionHeadFont}{.}{.5em}{}
\newtheoremstyle{fayorg-example}
  {0pt}{0pt}{\FayorgExampleBodyFont}{}
  {\FayorgExampleHeadFont}{.}{.5em}{}
\newtheoremstyle{fayorg-remark}
  {0pt}{0pt}{\FayorgRemarkBodyFont}{}
  {\FayorgRemarkHeadFont}{.}{.5em}{}
\newtheoremstyle{fayorg-exercise}
  {0pt}{0pt}{\FayorgExerciseBodyFont}{}
  {\FayorgExerciseHeadFont}{.}{.5em}{}
\newtheoremstyle{fayorg-explanation}
  {0pt}{0pt}{\FayorgExplanationBodyFont}{}
  {\FayorgExplanationHeadFont}{.}{.5em}{}
\newtheoremstyle{fayorg-answer}
  {0pt}{0pt}{\FayorgAnswerBodyFont}{}
  {\FayorgAnswerHeadFont}{.}{.5em}{}
\newtheoremstyle{fayorg-conjecture}
  {0pt}{0pt}{\FayorgConjectureBodyFont}{}
  {\FayorgConjectureHeadFont}{.}{.5em}{}

\theoremstyle{fayorg-theorem}
\newtheorem{theorem}{Theorem}[section]
\newtheorem{lemma}[theorem]{Lemma}
\newtheorem{proposition}[theorem]{Proposition}
\newtheorem{corollary}[theorem]{Corollary}

\theoremstyle{fayorg-definition}
\newtheorem{fayorgdefinition}[theorem]{Definition}
\newtheorem*{fayorgdefinitionstar}{Definition}

\theoremstyle{fayorg-example}
\newtheorem{example}[theorem]{Example}
\newtheorem*{eg}{Example}

\theoremstyle{fayorg-remark}
\newtheorem*{remark}{Remark}

\theoremstyle{fayorg-exercise}
\newtheorem{fayorgexercise}{Exercise}[section]

\theoremstyle{fayorg-explanation}
\newtheorem*{fayorgexplanation}{Proof}

\theoremstyle{fayorg-answer}
\newtheorem*{answer}{Answer}

\theoremstyle{fayorg-conjecture}
\newtheorem{conjecture}[theorem]{Conjecture}

% Definitions keep the legacy behavior of starting the body on a new line when
% an optional name is supplied.
\NewDocumentEnvironment{definition}{o}
  {
    \IfNoValueTF{#1}
      {\begin{fayorgdefinition}}
      {\begin{fayorgdefinition}[#1]\leavevmode\par\nobreak\smallskip}
  }
  {\end{fayorgdefinition}}

\NewDocumentEnvironment{definition*}{o}
  {
    \IfNoValueTF{#1}
      {\begin{fayorgdefinitionstar}}
      {\begin{fayorgdefinitionstar}[#1]\leavevmode\par\nobreak\smallskip}
  }
  {\end{fayorgdefinitionstar}}

% Exercise accepts an optional subtitle and an optional name. With one value it
% is shown in italics; with two values they are combined as in the old setup.
\NewDocumentEnvironment{exercise}{o o}
  {
    \IfNoValueTF{#1}
      {\begin{fayorgexercise}}
      {
        \IfNoValueTF{#2}
          {\begin{fayorgexercise}[\textit{#1}]\leavevmode\par\nobreak\smallskip}
          {\begin{fayorgexercise}[#2, \textit{#1}]\leavevmode\par\nobreak\smallskip}
      }
  }
  {\end{fayorgexercise}}

\NewDocumentCommand{\Answer}{}
  {
    \par\vspace{\baselineskip}
    {\FayorgExerciseHeadFont Answer:}\par\smallskip
  }

\NewDocumentEnvironment{explanation}{o}
  {
    \vspace{-10pt}\pushQED{\(\circledast\)}
    \IfNoValueTF{#1}
      {\begin{fayorgexplanation}}
      {\begin{fayorgexplanation}[#1]}
  }
  {\null\hfill\popQED\end{fayorgexplanation}}

\tcolorboxenvironment{theorem}{fayorg-theorem}
\tcolorboxenvironment{lemma}{fayorg-theorem}
\tcolorboxenvironment{proposition}{fayorg-theorem}
\tcolorboxenvironment{corollary}{fayorg-theorem}
\tcolorboxenvironment{fayorgdefinition}{fayorg-definition}
\tcolorboxenvironment{fayorgdefinitionstar}{fayorg-definition}
\tcolorboxenvironment{example}{fayorg-example}
\tcolorboxenvironment{eg}{fayorg-example}
\tcolorboxenvironment{remark}{fayorg-remark}
\tcolorboxenvironment{fayorgexercise}{fayorg-exercise}
\tcolorboxenvironment{fayorgexplanation}{fayorg-explanation}
\tcolorboxenvironment{answer}{fayorg-answer}
\tcolorboxenvironment{conjecture}{fayorg-conjecture}

\newtcolorbox{note}[1][]{fayorg-note,title={Note},#1}

% Options passed to code are listings options, for example
% \begin{code}[language=Python].
\newtcblisting{code}[1][]{
  fayorg-code,
  listing only,
  listing options={style=fayorg-code,#1}
}
